\chapter*{不定积分每日一题}

\refer{《积分的方法与技巧》金玉明\ 顾新身\ 毛瑞庭\ /\ 编著}

\exercise{求积分$\int{\frac{9-7x}{x^2+12x+38}\d x}$}%在花括号内输入题目

\teacher{张晓鹏}%在花括号内输入你的名字

\postdate{2023年4月4日}%在花括号内输入你想要发布那一天的日期

\result{-\frac{7}{2}\ln{\abs{x^2+12x+38}}+\frac{51}{\sqrt{2}}\arctan{\frac{x+6}{\sqrt{2}}}+C}%在花括号内输入本题结果

\analysis{
    关于这个的积分，我们可以归结为$\int{\frac{px+q}{ax^2+bx+c}}$这一类积分.而这个积分在有理分式积分中是也经常出现的，我们可以利用凑微分的操作对它进行化简，具体操作如下：
    \formula{
        \begin{aligned}
            \int{\frac{px+q}{ax^2+bx+c}\d x}&=\int{\frac{\dfrac{p}{2a}\qty(2ax+b)+\qty(q-\dfrac{pb}{2a})}{ax^2+bx+c}\d x}\\
            &=\frac{p}{2a}\int{\frac{2ax+b}{ax^2+bx+c}\d x}+\qty(q-\frac{pb}{2a})\int{\frac{\d x}{ax^2+bx+c}}\\
            &=\frac{p}{2a}\int{\frac{\d\qty(ax^2+bx+c)}{ax^2+bx+c}}+\qty(q-\frac{pb}{2a})\int{\frac{\d x}{ax^2+bx+c}}\\
            &=\frac{p}{2a}\ln{\abs{ax^2+bx+c}}+\qty(q-\frac{pb}{2a})\int{\frac{\d x}{ax^2+bx+c}}
        \end{aligned}
    }
    而剩下的\hyperlink{4.4}{$\int{\frac{\d x}{ax^2+bx+c}}$}我们昨天已经分三种情况讨论过啦，可以参考昨天的内容.
    
    对应今天的这个题目，为了方便可以直接代入（实际做题推一下即可，还是很方便的，分子凑分母的导数即可）
    \formula{\int{\frac{9-7x}{x^2+12x+38}\d x}=-\frac{7}{2}\ln{\abs{x^2+12x+38}}+\qty[9-\frac{-7\times 12}{2}]\int{\frac{\d x}{x^2+12x+38}}}
    而对于$\int{\frac{\d x}{x^2+12x+38}}$，属于$b^2=144<4ac=152$的情况，那么
    \formula{
        \begin{aligned}
        \int{\frac{\d x}{x^2+12x+38}}&=\frac{2}{\sqrt{152-144}}\arctan{\frac{2x+12}{\sqrt{152-144}}}+C\\
        &=\frac{1}{\sqrt{2}}\arctan{\frac{x+6}{\sqrt{2}}}+C
        \end{aligned}
        }
        因此可以得到最后答案（答案并不重要，思想重要！）
            \formula{\int{\frac{9-7x}{x^2+12x+38}\d x}=-\frac{7}{2}\ln{\abs{x^2+12x+38}}+\frac{51}{\sqrt{2}}\arctan{\frac{x+6}{\sqrt{2}}}+C}
}%在花括号内输入分析